\documentclass[10pt]{article}
% Shared typographic and algebraic source.
\usepackage[letterpaper,margin=0.86in]{geometry}
\usepackage[T1]{fontenc}
\usepackage[utf8]{inputenc}
\usepackage{lmodern}
\usepackage{microtype}
\usepackage{amsmath,amssymb,amsthm,mathtools,mathrsfs}
\usepackage{booktabs,array,tabularx,longtable}
\usepackage{enumitem}
\usepackage[dvipsnames]{xcolor}
\usepackage{fancyhdr,titlesec}
\usepackage[hidelinks]{hyperref}

\definecolor{ink}{HTML}{1F2A33}
\definecolor{teal}{HTML}{226A73}
\definecolor{muted}{HTML}{65747E}
\definecolor{line}{HTML}{C5D1D3}
\definecolor{wash}{HTML}{F3F7F6}

\setlength{\parindent}{0pt}
\setlength{\parskip}{0.46em}
\setlength{\emergencystretch}{2em}
\allowdisplaybreaks[2]
\renewcommand{\arraystretch}{1.16}
\setlist{nosep,leftmargin=1.5em}

\pagestyle{fancy}
\fancyhf{}
\fancyhead[L]{\footnotesize\sffamily\color{muted}COPPALINI ARCHITECTURE}
\fancyhead[R]{\footnotesize\sffamily\color{muted}PHYSICAL REALITY INSIDE CONSTRAINT}
\fancyfoot[C]{\footnotesize\sffamily\color{muted}\thepage}
\renewcommand{\headrulewidth}{0.35pt}
\renewcommand{\headrule}{\hbox to\headwidth{\color{line}\leaders\hrule height \headrulewidth\hfill}}
\setlength{\headheight}{19pt}

\titleformat{\section}{\Large\bfseries\color{ink}}{\thesection.}{0.55em}{}
\titleformat{\subsection}{\large\bfseries\color{teal}}{\thesubsection.}{0.5em}{}
\titlespacing*{\section}{0pt}{1.42em}{0.48em}
\titlespacing*{\subsection}{0pt}{1.02em}{0.32em}

\newtheorem{theorem}{Theorem}
\newtheorem{proposition}[theorem]{Proposition}
\newtheorem{lemma}[theorem]{Lemma}
\theoremstyle{definition}
\newtheorem{definition}[theorem]{Definition}
\theoremstyle{remark}
\newtheorem{remark}[theorem]{Boundary}

\newcommand{\RR}{\mathbb R}
\newcommand{\CC}{\mathbb C}
\newcommand{\QQ}{\mathbb Q}
\newcommand{\ZZ}{\mathbb Z}
\newcommand{\Id}{I}
\newcommand{\tr}{\operatorname{tr}}
\newcommand{\im}{\operatorname{im}}
\newcommand{\kerop}{\operatorname{ker}}
\newcommand{\Span}{\operatorname{span}}
\newcommand{\diag}{\operatorname{diag}}
\newcommand{\snoc}{\operatorname{snoc}}
\newcommand{\Admit}{\operatorname{Admit}}
\newcommand{\Gate}{\operatorname{Gate}}
\newcommand{\Typed}{\operatorname{Typed}}
\newcommand{\Future}{\operatorname{future}}
\newcommand{\Scope}{\operatorname{sc}}
\newcommand{\SilScope}{\mathsf{Silence}_{\rm scope}}
\newcommand{\SilFit}{\mathsf{Silence}_{\rm fit}}

\newcommand{\papertitle}[2]{%
  \begin{center}
  {\color{teal}\rule{\textwidth}{1.35pt}}\\[0.38in]
  {\footnotesize\sffamily\bfseries\color{teal} COPPALINI ARCHITECTURE}\\[0.16in]
  {\fontsize{31}{36}\selectfont\bfseries\color{ink}#1\par}
  \vspace{0.15in}
  {\Large\color{teal}#2\par}
  \vspace{0.22in}
  {\color{line}\rule{\textwidth}{0.55pt}}\\[0.12in]
  {\normalsize Anthony Vito Coppa}\par
  {\small 9 November 2026}\par
  \vspace{0.30in}
  \end{center}
}

\newcommand{\sharedroot}{%
Let
\[
Q=\begin{pmatrix}1&1\\1&0\end{pmatrix},\qquad
\Omega=\begin{pmatrix}0&1\\-1&0\end{pmatrix},\qquad
C=-\Omega,
\]
\[
G=[Q,\Omega]=\begin{pmatrix}-2&1\\1&2\end{pmatrix},\qquad
H=\frac{G}{\sqrt5},\qquad
S=\frac{2Q-I}{\sqrt5}.
\]
The common exact identities are
\[
Q^2=Q+I,\qquad C^2=-I,\qquad G^2=5I,\qquad H^2=S^2=I,
\]
\[
H=CS,\qquad SG=\sqrt5\Omega=-GS,\qquad
L_s:=e^{sS},\qquad
L_s^{\mathsf T}GL_s=G,\qquad HL_sH=L_{-s}.
\]
}

\newcommand{\sharedoccurrence}[2]{%
A physical field has an interface
\[
\begin{aligned}
\mathfrak I_J:=(&J,\Gamma_J,\operatorname{Src}_J;
A_J,X_J,Y_J,T_J,\Sigma_J,\operatorname{Con}_J,\\
&\rho_J,O_J,K_J,B_J,\theta_J),
\qquad
\rho_J:X_J\to Y_J,\quad
O_J:X_J\to\mathcal O_J,\quad
\theta_J:X_J\to\{0,1\}.
\end{aligned}
\tag{#1}
\]
Its occurrence type is
\[
\begin{aligned}
\mathsf{Occ}_J:=\{&
(\mathfrak I_J;\Theta,x,\rho_J(x),\kappa,
\boldsymbol\varepsilon,\theta_J(x)):\\
&\Theta\in A_J,\ x\in X_J,\ \kappa\in K_J,\
\boldsymbol\varepsilon\in\RR_{\ge0}^{3}\}.
\end{aligned}
\tag{#2}
\]
An indexed occurrence $\mathsf{Occ}_{J,n}$ is an element of
$\mathsf{Occ}_J$.
}

\newcommand{\sharedmatching}{%
Let
\[
\mathcal A_8:=\{\infty\}\sqcup\ZZ_7,
\qquad
(\infty,0,1,2,3,4,5,6)
=(\mathsf{Gr},\mathsf{Li},\mathsf{Me},\mathsf{Ti},
\mathsf{Hz},\mathsf{Si},\mathsf{Sc},\mathsf{Mo}).
\]
For $r\in\ZZ_7$, set
\[
\mathsf M_r=\{\{\infty,r\}\}\cup
\{\{r+j,r-j\}:j=1,2,3\}.
\]
Then $E(K_8)=\bigsqcup_{r\in\ZZ_7}\mathsf M_r$, and the canonical matching is
\[
\mathsf M_0=\{\mathsf{Gravity\! -\! Light},\;
\mathsf{Memory\! -\! Motion},\;
\mathsf{Time\! -\! Scope},\;
\mathsf{Horizon\! -\! Silence}\}.
\]
}

\hypersetup{pdftitle={Suture Field Dynamics: Regular Hamiltonian Reduction},pdfauthor={Anthony Vito Coppa},pdfsubject={Constrained action, regular suture reduction, and Einstein-Maxwell degrees of freedom}}
\newcommand{\hateta}{\widehat\eta}
\newcommand{\hatC}{\widehat C}
\newcommand{\Fsol}{\mathcal F_{\rm sol}}
\newcommand{\rankop}{\operatorname{rank}}
\newcommand{\DoF}{\operatorname{DoF}}
\begin{document}
\papertitle{Suture Field Dynamics}{Regular Hamiltonian reduction}

\begin{abstract}
The suture field supplies an eight-dimensional internal carrier for a
Lorentzian metric and an antisymmetric flux form. This note gives the exact
condition under which its multiplier action has the same propagating content
as Einstein--Maxwell theory. The condition is the rank-sixteen regularity of
the pointwise suture map from the thirty-two coframe coordinates to the ten
metric and six flux coordinates. On that branch the remaining sixteen coframe
coordinates are vertical gauge coordinates; the multiplier sector is
second-class; and the reduced Hamiltonian carries two gravitational and two
electromagnetic polarizations.
\end{abstract}

\section{Constrained action}

Let $\Pi^A{}_{\mu}$ be a rank-four $E$-valued coframe, with $A=1,\ldots,8$.
Let $B^{\mu\nu}=-B^{\nu\mu}$ and
$\lambda^{\mu\nu}=\lambda^{\nu\mu}$. Fix a nonzero constant calibration
$\chi$. Define
\[
h_{\mu\nu}(\Pi):=\hateta_{AB}\Pi^A{}_{\mu}\Pi^B{}_{\nu},
\qquad
f_{\mu\nu}(\Pi):=\chi\hatC_{AB}\Pi^A{}_{\mu}\Pi^B{}_{\nu}.
\tag{SD.0}
\]
The action is
\[
\begin{aligned}
\mathcal S=\int_M\Bigg[&\frac{c^3}{16\pi G_N}(R(g)-2\Lambda)
-\frac1{4\mu_0}F_{\mu\nu}F^{\mu\nu}\\
&+\frac12B^{\mu\nu}\bigl(F_{\mu\nu}-f_{\mu\nu}(\Pi)\bigr)
+\lambda^{\mu\nu}\bigl(g_{\mu\nu}-h_{\mu\nu}(\Pi)\bigr)\Bigg]
\sqrt{-g}\,d^4x,
\qquad F_{\mu\nu}=2\partial_{[\mu}A_{\nu]}.
\end{aligned}
\tag{SD.1}
\]
Variation of the multipliers gives the two suture conditions
\[
F_{\mu\nu}=f_{\mu\nu}(\Pi),
\qquad g_{\mu\nu}=h_{\mu\nu}(\Pi).
\tag{SD.2}
\]
They are constraint equations. Their multiplier momenta belong to the
canonical system and remain present until Dirac reduction.
In particular, variation with respect to $B$ and $lambda$ gives SD.2;
it does not eliminate $B$ or $lambda$. Their values can only be related to
the other equations after the coupled Euler system has been imposed, and
that later relation does not remove their primary momenta from the Dirac
analysis.

\section{The pointwise suture map}

At one spacetime address, define
\[
\mathscr S:\operatorname{Hom}(\RR^4,\RR^8)_{\mathrm{rank}\,4}
\longrightarrow \operatorname{Sym}^2(\RR^4)^*\oplus\Lambda^2(\RR^4)^*,
\qquad
\Pi\longmapsto(h(\Pi),f(\Pi)).
\tag{SD.3}
\]
The domain has dimension $32$ and the target has dimension $10+6=16$. Its
differential is
\[
\begin{aligned}
(D\mathscr S_\Pi\,\delta\Pi)^{(g)}_{\mu\nu}
&=\hateta_{AB}\bigl(\delta\Pi^A{}_{\mu}\Pi^B{}_{\nu}
+\Pi^A{}_{\mu}\delta\Pi^B{}_{\nu}\bigr),\\
(D\mathscr S_\Pi\,\delta\Pi)^{(F)}_{\mu\nu}
&=\chi\hatC_{AB}\bigl(\delta\Pi^A{}_{\mu}\Pi^B{}_{\nu}
+\Pi^A{}_{\mu}\delta\Pi^B{}_{\nu}\bigr).
\end{aligned}
\tag{SD.4}
\]

\begin{definition}[Regular suture branch]
For the fixed $\chi\ne0$, the regular branch is the open locus on which
\[
\rankop D\mathscr S_\Pi=16.
\tag{SD.5}
\]
\end{definition}

An exact certificate that this locus is nonempty is
\[
\Pi_\star=
\begin{pmatrix}
-1&-2& 1& 0\\
 2&-1& 1& 0\\
-1&-2& 1&-2\\
-1& 1&-1& 0\\
-1& 0& 0&-1\\
 0& 1& 1&-1\\
 2& 2& 1&-2\\
-1&-1& 0&-1
\end{pmatrix}.
\tag{SD.6}
\]
For $\chi=1$, exact row reduction gives
\[
\rankop D\mathscr S_{\Pi_\star}=16,
\qquad
h(\Pi_\star)=
\begin{pmatrix}
10&-5& 9& 0\\
-5& 9&-1& 0\\
 9&-1& 6&-5\\
 0& 0&-5&11
\end{pmatrix}.
\tag{SD.7}
\]
Its leading principal minors are
\[
10,\qquad65,\qquad-259,\qquad-4474,
\tag{SD.8}
\]
so the $LDL^{\mathsf T}$ pivot signs are $+,+,-,+$ and
$h(\Pi_\star)$ has Lorentzian signature $(3,1)$. The determinant of a
nonzero $16\times16$ minor of $D\mathscr S$ remains nonzero on a
neighborhood of $\Pi_\star$, which establishes the stated open regular
branch.

\section{Internal vertical directions}

The linear internal stabilizer is
\[
\mathfrak k=
\left\{X\in\mathfrak{gl}(8,\RR):
X^{\mathsf T}\hateta+\hateta X=0,\quad
X^{\mathsf T}\hatC+\hatC X=0\right\}.
\tag{SD.9}
\]
Exact linear reduction gives
\[
\dim\mathfrak k=10,
\qquad
\mathfrak k\cong\mathfrak{so}(1,1)\oplus\mathfrak u(3).
\tag{SD.10}
\]
The ten linear directions $\delta\Pi=X\Pi$ preserve both $h$ and $f$.
On the regular branch,
\[
\dim\ker D\mathscr S_\Pi=32-16=16.
\tag{SD.11}
\]
The other six vertical directions are field-dependent coordinates along the
same local fiber of $\mathscr S$. They carry no metric or flux coordinate.

\section{The coframe bracket matrix}

The primary coframe momenta are $P_A{}^\mu\approx0$. Their stabilization
contains
\[
\Psi_A{}^\mu=\sqrt{-g}\left(
2\lambda^{\mu\nu}\hateta_{AB}\Pi^B{}_{\nu}
+\chi B^{\mu\nu}\hatC_{AB}\Pi^B{}_{\nu}\right)\approx0.
\tag{SD.12}
\]
With the canonical convention
$\{\Pi^A{}_{\mu}(x),P_B{}^\nu(y)\}=\delta^A_B\delta_\mu^\nu
\delta^3(x-y)$, the cross-bracket is
\[
\{P_A{}^\mu(x),\Psi_B{}^\nu(y)\}
=-\Delta_{AB}{}^{\mu\nu}(x)\,\delta^3(x-y),
\qquad
\Delta_{AB}{}^{\mu\nu}
=2\sqrt{-g}\left(
\lambda^{\mu\nu}\hateta_{AB}
+\frac\chi2B^{\mu\nu}\hatC_{AB}\right).
\tag{SD.13}
\]
Changing the Poisson-bracket ordering reverses this displayed overall sign;
the rank statement is unchanged.
At the vacuum multiplier point
\[
\lambda^{\mu\nu}=0,\qquad B^{\mu\nu}=0,
\tag{SD.14}
\]
one has $\Delta=0$. Thus nondegeneracy of $\hateta$ and $\hatC$, and rank
four of $\Pi$, do not establish invertibility of $\Delta$. The $32$ coframe
primaries and $32$ displayed stabilizations therefore do not form a globally
invertible second-class block.

This point occurs on a Lorentzian suture vacuum. Let the four columns of
$\Pi_{\rm vac}$ be the first coordinate vector of each $2\times2$ carrier
block. Then
\[
h(\Pi_{\rm vac})=\diag(-2,1,1,1),
\qquad
f(\Pi_{\rm vac})=0.
\tag{SD.14a}
\]
For $\Lambda=0$, the constant metric $g=h(\Pi_{\rm vac})$ and $A=0$ satisfy
the vacuum Einstein--Maxwell equations with the multiplier values in SD.14.

\section{Regular multiplier reduction}

For an unambiguous Legendre transform, introduce an independent antisymmetric
field $\mathcal F_{\mu\nu}$ and an antisymmetric multiplier $Y^{\mu\nu}$:
\[
\begin{aligned}
\mathcal S_{\rm fo}=\int_M\Bigg[&\frac{c^3}{16\pi G_N}(R-2\Lambda)
-\frac1{4\mu_0}\mathcal F_{\mu\nu}\mathcal F^{\mu\nu}\\
&+\frac12Y^{\mu\nu}\bigl(\mathcal F_{\mu\nu}-2\partial_{[\mu}A_{\nu]}\bigr)
+\frac12B^{\mu\nu}\bigl(\mathcal F_{\mu\nu}-f_{\mu\nu}(\Pi)\bigr)\\
&+\lambda^{\mu\nu}\bigl(g_{\mu\nu}-h_{\mu\nu}(\Pi)\bigr)\Bigg]
\sqrt{-g}\,d^4x.
\end{aligned}
\tag{SD.15}
\]
Variation of $Y$ returns $\mathcal F=dA$, so SD.15 is a first-order form of
SD.1. Variation of $\mathcal F$ gives
\[
Y^{\mu\nu}+B^{\mu\nu}=\frac{1}{\mu_0}\mathcal F^{\mu\nu}.
\tag{SD.15a}
\]
Thus $B$ remains an auxiliary multiplier variable through the canonical
reduction; SD.15a is not a license to remove its primary constraint before
that reduction.

By the submersion theorem, local coordinates on the regular coframe carrier
are
\[
\Pi\longleftrightarrow(q^I,\xi^a),
\qquad
q^I=\mathscr S^I(\Pi),
\qquad I,a=1,\ldots,16.
\tag{SD.16}
\]
In these coordinates, the multiplier sector has normal form
\[
\mathcal L_{\rm mult}=\sqrt{-g}\;m_I\bigl(X^I-q^I\bigr),
\qquad
X=(g,\mathcal F),
\qquad
m=(\lambda,B/2).
\tag{SD.17}
\]
Its three Euler relations are
\[
X^I-q^I=0,
\qquad m_I=0,
\qquad \frac{\delta\mathcal L_{\rm mult}}{\delta\xi^a}=0.
\tag{SD.18}
\]
The $\xi^a$ are vertical gauge coordinates. The remaining auxiliary
constraints admit the local Darboux pairing
\[
(P_{qI},q^I-X^I),
\qquad
(P_{mI},m_I),
\tag{SD.19}
\]
which gives $64$ second-class constraints. The $16$ constraints
\[
P_{\xi a}\approx0
\tag{SD.20}
\]
are first class and generate vertical suture motions.

More explicitly, write
\[
\mathcal C_I:=q^I-X^I.
\tag{SD.20a}
\]
\par\noindent
The auxiliary Hamiltonian has the local form
\[
H_{\rm aux}=\int d^3x\,
\bigl(m_I\mathcal C_I
+u_m^IP_{mI}+u_q^IP_{qI}+u_\xi^aP_{\xi a}\bigr).
\tag{SD.20b}
\]
\par\noindent
Its primary stabilizations are
\[
\dot P_{mI}\approx0\Longrightarrow\mathcal C_I\approx0,
\qquad
\dot P_{qI}\approx0\Longrightarrow m_I\approx0,
\qquad
\dot P_{\xi a}\approx0\Longrightarrow0.
\tag{SD.20c}
\]
\par\noindent
With ordered constraint vector
\[
\boldsymbol\zeta=(P_q,\mathcal C;P_m,m),
\tag{SD.20d}
\]
\par\noindent
the second-class matrix is
\[
\{\boldsymbol\zeta(x),\boldsymbol\zeta(y)\}
=
\begin{pmatrix}
0&-I_{16}&0&0\\
I_{16}&0&0&0\\
0&0&0&-I_{16}\\
0&0&I_{16}&0
\end{pmatrix}\delta^3(x-y).
\tag{SD.20e}
\]
\par\noindent
It is invertible. Stabilization of $\mathcal C_I$ and $m_I$ fixes
$u_q^I$ and $u_m^I$; it creates no tertiary auxiliary constraint.

\begin{theorem}[Regular suture reduction]
On the regular branch, Dirac reduction of SD.15 is locally equivalent to
\[
\mathcal S_{\rm red}=\int_M\left[
\frac{c^3}{16\pi G_N}(R-2\Lambda)
-\frac1{4\mu_0}F_{\mu\nu}F^{\mu\nu}
\right]\sqrt{-g}\,d^4x,
\qquad F=dA.
\tag{SD.21}
\]
\par\noindent
\end{theorem}

\begin{proof}
SD.18 removes $q$ and $m$ while SD.20 removes the vertical coordinates.
Variation of $Y$ returns $\mathcal F=dA$. Substitution leaves SD.21.
\end{proof}

\section{ADM Hamiltonian and stabilization}

Let
\[
K_{ij}=\frac1{2N}\bigl(\dot\gamma_{ij}-D_iN_j-D_jN_i\bigr),
\qquad
\pi^{ij}=\frac{c^3\sqrt\gamma}{16\pi G_N}
\bigl(K^{ij}-\gamma^{ij}K\bigr),
\tag{SD.22}
\]
\par\noindent
and, after the multiplier reduction, let
\[
E^i=-\frac{\sqrt{-g}}{\mu_0}F^{0i}.
\tag{SD.23}
\]
\par\noindent
The reduced primary constraints are
\[
\pi_N\approx0,\qquad \pi_i\approx0,\qquad E^0\approx0.
\tag{SD.24}
\]
\par\noindent
The total Hamiltonian is
\[
H_T=\int d^3x\left(
N\mathcal H_\perp+N^i\mathcal H_i+A_0\mathcal G
+u_N\pi_N+u^i\pi_i+u_0E^0\right),
\tag{SD.25}
\]
\par\noindent
where
\[
\begin{aligned}
\mathcal H_\perp={}&
\frac{16\pi G_N}{c^3\sqrt\gamma}
\left(\pi^{ij}\pi_{ij}-\frac12\pi^2\right)
-\frac{c^3\sqrt\gamma}{16\pi G_N}\bigl({}^{(3)}R-2\Lambda\bigr)\\
&+\frac{\mu_0}{2\sqrt\gamma}\gamma_{ij}E^iE^j
+\frac{\sqrt\gamma}{4\mu_0}F_{ij}F^{ij},\\
\mathcal H_i={}&-2\gamma_{ik}D_j\pi^{jk}+F_{ij}E^j,
\qquad
\mathcal G=-\partial_iE^i.
\end{aligned}
\tag{SD.26}
\]
\par\noindent
Preservation of SD.24 gives
\[
\dot\pi_N\approx0\Longrightarrow\mathcal H_\perp\approx0,
\qquad
\dot\pi_i\approx0\Longrightarrow\mathcal H_i\approx0,
\qquad
\dot E^0\approx0\Longrightarrow\mathcal G\approx0.
\tag{SD.27}
\]
\par\noindent
The scalar, momentum, and Gauss constraints close under the standard
hypersurface-deformation and $U(1)$ gauge algebra. Their stabilizations fix
only the corresponding Dirac multipliers and produce no tertiary physical
constraint.

\section{Degree count}

On the regular Darboux chart of the direct action SD.1, the phase-space
dimensions per spatial address are:
\[
\begin{array}{c|c}
\text{sector}&\text{phase-space dimension}\\ \hline
(\gamma_{ij},N,N^i) &20\\
A_\mu &8\\
\Pi^A{}_{\mu}&64\\
B^{\mu\nu}&12\\
\lambda^{\mu\nu}&20\\ \hline
\text{total}&124
\end{array}
\tag{SD.28}
\]
There are ten Einstein--Maxwell first-class constraints and sixteen vertical
suture first-class constraints. The regular multiplier sector supplies sixty
four second-class constraints. Therefore
\[
\DoF=
\frac12\bigl(124-2(10+16)-64\bigr)=4.
\tag{SD.29}
\]
Equivalently, after exact reduction the Einstein--Maxwell phase space has
dimension $28$ and ten first-class constraints:
\[
\DoF_{\rm red}=\frac12\bigl(28-2\cdot10\bigr)=4.
\tag{SD.30}
\]
The four propagating degrees of freedom are the two gravitational and two
electromagnetic transverse polarizations.

\section{Discrete carrier}

The $\ZZ_7$ Fourier action of the suture carrier preserves $\hateta$ and
$\hatC$. In SD.1 it is an algebraic fiber symmetry and contributes no new
canonical pair. A local connection or a measured holonomy coupling enters
only through an additional specified dynamical term and calibrated
occurrence record.

\end{document}
